% !TeX program = pdflatex
\documentclass[11pt,a4paper]{amsart}

% Shared preamble for course notes and exercise sets.

\usepackage[T1]{fontenc}
\usepackage{lmodern}
\usepackage[english]{babel}

\usepackage[a4paper,margin=30mm]{geometry}
\usepackage{microtype}
\usepackage{graphicx}
\usepackage{enumitem}

\usepackage{amsmath}
\usepackage{amssymb}
\usepackage{amsthm}
\usepackage{mathtools}
\usepackage{aliascnt}

\usepackage[hidelinks,pdfusetitle]{hyperref}

\numberwithin{equation}{section}
\setlist{itemsep=0.25em,topsep=0.5em}

\theoremstyle{plain}
\newtheorem{theorem}{Theorem}[section]
\newaliascnt{lemma}{theorem}
\newtheorem{lemma}[lemma]{Lemma}
\aliascntresetthe{lemma}
\newaliascnt{proposition}{theorem}
\newtheorem{proposition}[proposition]{Proposition}
\aliascntresetthe{proposition}
\newaliascnt{corollary}{theorem}
\newtheorem{corollary}[corollary]{Corollary}
\aliascntresetthe{corollary}

\theoremstyle{definition}
\newaliascnt{definition}{theorem}
\newtheorem{definition}[definition]{Definition}
\aliascntresetthe{definition}
\newaliascnt{example}{theorem}
\newtheorem{example}[example]{Example}
\aliascntresetthe{example}
\newaliascnt{problem}{theorem}
\newtheorem{problem}[problem]{Problem}
\aliascntresetthe{problem}

\theoremstyle{remark}
\newaliascnt{remark}{theorem}
\newtheorem{remark}[remark]{Remark}
\aliascntresetthe{remark}

\usepackage[nameinlink,noabbrev]{cleveref}

\crefname{theorem}{theorem}{theorems}
\crefname{lemma}{lemma}{lemmas}
\crefname{proposition}{proposition}{propositions}
\crefname{corollary}{corollary}{corollaries}
\crefname{definition}{definition}{definitions}
\crefname{example}{example}{examples}
\crefname{problem}{problem}{problems}
\crefname{remark}{remark}{remarks}

\DeclareMathOperator{\im}{im}
\DeclareMathOperator{\rank}{rank}
\DeclarePairedDelimiter{\abs}{\lvert}{\rvert}
\DeclarePairedDelimiter{\norm}{\lVert}{\rVert}
\DeclarePairedDelimiter{\set}{\lbrace}{\rbrace}


\title{Analyysi I: Limits}
\author{Joona Piirainen}
\date{Autumn 2026}

\begin{document}

\maketitle

\begin{abstract}
  Notes for the University of Helsinki course MAT11003, taught in Finnish
  from 01.08.2026--17.10.2026.
\end{abstract}

\tableofcontents

\section{Course information}

The official course information is available on the
\href{https://studies.helsinki.fi/kurssit/toteutus/hy-opt-cur-2627-0f6de353-53e0-4057-8e8e-b45cab4f2a8d}{University
of Helsinki course page}.

\section{Lecture notes}

\subsection{01.09.2026}

\subsubsection{The absolute value function}

\begin{definition}\label{def:abs}
  For any real number \(x \in \RR\), the \emph{absolute value} (or
  \emph{modulus}) of \(x\), denoted by \(\abs{x}\), is defined by
  \[
    \abs{x} =
    \begin{cases}
      x,  & \text{if } x \geq 0, \\
      -x, & \text{if } x < 0.
    \end{cases}
  \]
\end{definition}

Geometrically, \(\abs{x-y}\) is the distance between the points \(x\) and \(y\)
on the real line.

\begin{example}
  Fix \(y \in \RR\) and let \(h>0\). The inequality \(\abs{x-y}<h\) can be
  written in several equivalent ways:
  \begin{align*}
    \abs{x-y}<h
    &\iff -h<x-y<h \\
    &\iff y-h<x<y+h \\
    &\iff x\in(y-h,y+h).
  \end{align*}
  Thus, \(\abs{x-y}<h\) means that \(x\) lies in the open interval of radius
  \(h\) centred at \(y\).
\end{example}

\begin{example}
  As a concrete instance of the preceding equivalences,
  \begin{align*}
    \abs{2x-1}<3
    &\iff -3<2x-1<3 \\
    &\iff -2<2x<4 \\
    &\iff -1<x<2 \\
    &\iff x\in(-1,2).
  \end{align*}
\end{example}

\begin{theorem}[Triangle and reverse triangle inequalities]
  For all \(x,y\in\RR\),
  \begin{align*}
    \abs{x+y} &\leq \abs{x}+\abs{y}, \\
    \abs{\abs{x}-\abs{y}} &\leq \abs{x-y}.
  \end{align*}
\end{theorem}

\begin{proof}
  By \cref{def:abs},
  \begin{align*}
    -\abs{x}&\leq x\leq\abs{x}, \\
    -\abs{y}&\leq y\leq\abs{y}.
  \end{align*}
  Adding the inequalities gives
  \[
    -\bigl(\abs{x}+\abs{y}\bigr)
    \leq x+y
    \leq \abs{x}+\abs{y},
  \]
  which is equivalent to \(\abs{x+y}\leq\abs{x}+\abs{y}\).

  Applying this triangle inequality to \(x=(x-y)+y\) gives
  \[
    \abs{x}-\abs{y}\leq\abs{x-y}.
  \]
  Interchanging \(x\) and \(y\) similarly gives
  \(\abs{y}-\abs{x}\leq\abs{x-y}\). Therefore
  \[
    -\abs{x-y}\leq\abs{x}-\abs{y}\leq\abs{x-y},
  \]
  which is equivalent to the reverse triangle inequality.
\end{proof}

\subsection{03.09.2026}

\subsubsection{More on absolute values}

\begin{theorem}
  We have
  \[
    \abs{xy} = \abs{x}\abs{y}
  \]
  for all \(x, y \in \RR\).
\end{theorem}

\begin{proof}
  We prove by splitting into four cases. Firstly, if \(x,y \geq 0\), we have
  \(xy \geq 0\), so \(\abs{xy} = xy\). But in this case also both
  \(\abs{x} = x\) and \(\abs{y} = y\), so \(\abs{xy} = xy = \abs{x}\abs{y}\).
  Secondly, if \(x, y < 0\), we have \(xy > 0\),
  \(\abs{xy} = xy = (-x)(-y) = \abs{x}\abs{y}\). Now suppose
  \(x < 0\) and \(y \geq 0\). Then \(xy\leq 0\), and hence
  \(\abs{xy}=-xy=(-x)y=\abs{x}\abs{y}\). The
  case \(x \geq 0\), \(y < 0\) follows from commutativity of multiplication
  of real numbers.
\end{proof}

\begin{example}
  Let \(x \in (1, 3)\) and \(y \in (4, 8)\). Find an upper bound for
  \[\abs{x + y - 8}.\] Kolmioepäyhtälön nojalla saadaan
  \[
    \begin{aligned}
      \abs{x + y - 8} &= \abs{(x - 2) + (y - 6)} \\
      &\leq \abs{x - 2} + \abs{y - 6} \\
      &< 1 + 2 \\
      &= 3.
    \end{aligned}
  \]
  Thus we have found that \(3\) is an upper bound for
  \(\abs{x + y - 8}\).
\end{example}

\begin{example}
  Let \(x \in (2, 6)\) and \(y \in (2, 8)\). Find an upper boud for
  \[\abs{xy - 20}.\]
  We have \(\abs{x - 4} < 2\) and \(\abs{y - 5} < 3\). We have
  \[
    \begin{aligned}
      \abs{xy - 20} &= \abs{xy - 4y + 4y - 20} \\
      &\leq \abs{xy - 4y} + \abs{4y - 20} \\
      &= \abs{y}\abs{x - 4} + 4\abs{y - 5} \\
      &\leq 8 \cdot 2 + 4 \cdot 3 \\
      &= 28.
    \end{aligned}
  \]
  It follows that \(28\) is an upper bound for \(\abs{xy - 20}\).
\end{example}

\begin{example}
  Find a positive real number \(m\) such that
  \[\abs{x^2 - 16} \leq m\abs{x - 4},\]
  where \(x \in (0, 10)\).

  When \(x = 4\), we get \(0 \leq 0\), which holds with every \(m \in \RR\).
  Now suppose \(x \neq 4\). We simplify the inequality and obtain
  \[\abs{x + 4} \leq m.\]
  Since \(x \in (0, 10)\), \(\abs{x + 4} \in (4, 14)\). We have thus obtained
  \[\abs{x + 4} \le 14,\]
  so we can choose \(m = 14\).
\end{example}

\subsubsection{Field and ordering properties of \texorpdfstring{\(\RR\)}{R}}

\begin{definition}
  \((F, +, \cdot)\) is a field if and only if
  \begin{enumerate}
    \item The algebraic structure \((F, +)\) is an \emph{abelian group}.
    \item The algebraic structure \((F^{*}, \cdot)\) is an
      \emph{abelian group}, where \(F^{*} = F \setminus \{0_F\}\).
    \item The operation \(\cdot\) distributes over \(+\).
  \end{enumerate}
\end{definition}

\begin{theorem}
  \(\RR\) is a field.
\end{theorem}

\begin{proof}
  Consider the usual addition and multiplication on \(\RR\).

  First, \((\RR,+)\) is an abelian group. Addition is closed, associative,
  and commutative; \(0\) is an additive identity; and every \(x\in\RR\) has
  the additive inverse \(-x\).

  Next, \((\RR\setminus\{0\},\cdot)\) is an abelian group. Multiplication is
  associative and commutative, and \(1\) is its identity element. If
  \(x\neq 0\), then \(x^{-1}=1/x\) is a nonzero real number satisfying
  \[
    x\cdot x^{-1}=1.
  \]
  The product of two nonzero real numbers is nonzero, so multiplication is
  closed on \(\RR\setminus\{0\}\).

  Finally, multiplication distributes over addition:
  \[
    x(y+z)=xy+xz
  \]
  for all \(x,y,z\in\RR\). Therefore \((\RR,+,\cdot)\) is a field.
\end{proof}

\begin{definition}[Ordered field]
  An \emph{ordered field} is a field \((F,+,\cdot)\) equipped with a relation
  \(<\) such that, for all \(x,y,z\in F\),
  \begin{enumerate}
    \item exactly one of \(x<y\), \(x=y\), or \(x>y\) holds;
    \item if \(x<y\) and \(y<z\), then \(x<z\);
    \item if \(x<y\), then \(x+z<y+z\); and
    \item if \(0_F<x\) and \(0_F<y\), then \(0_F<xy\).
  \end{enumerate}
  Here \(x>y\) means \(y<x\). The first two conditions state that \(<\) is a
  strict total order, while the last two state that the order is compatible
  with the field operations.
\end{definition}

\begin{theorem}
  The field \((\RR,+,\cdot)\), equipped with the usual order \(<\), is an
  ordered field.
\end{theorem}

\begin{proof}
  We have already shown that \((\RR,+,\cdot)\) is a field. The usual strict
  order on \(\RR\) satisfies trichotomy and transitivity, so it satisfies the
  first two ordering axioms.

  Suppose that \(x<y\). Then \(0<y-x\), and for every \(z\in\RR\),
  \[
    (y+z)-(x+z)=y-x>0.
  \]
  Hence \(x+z<y+z\). Furthermore, if \(0<x\) and \(0<y\), then the usual
  sign rule for multiplication gives \(0<xy\). Thus the order
  is compatible with addition and multiplication, and \(\RR\) is an ordered
  field.
\end{proof}

\begin{theorem}[Uniqueness of additive identity]
  The identity element \(0_F\) of the abelian group \((F, +)\) is unique.
\end{theorem}
\begin{proof}
  Suppose that \(0_F\) and \(0'_F\) are both additive identity elements of
  \((F,+)\). Since \(0_F\) is an additive identity,
  \[
    0_F+0'_F=0'_F.
  \]
  Since \(0'_F\) is also an additive identity,
  \[
    0_F+0'_F=0_F.
  \]
  Therefore \(0_F=0'_F\).
\end{proof}

\begin{theorem}[Cancellation law]
  Let \(x,y,z \in \RR\). If \(x+y=x+z\), then \(y=z\).
\end{theorem}

\begin{proof}
  Using the additive identity, additive inverse, associativity, and the
  assumption \(x+y=x+z\), we obtain
  \begin{align*}
    y
    &=0+y \\
    &=((-x)+x)+y \\
    &=(-x)+(x+y) \\
    &=(-x)+(x+z) \\
    &=((-x)+x)+z \\
    &=0+z \\
    &=z.
  \end{align*}
\end{proof}

\subsection{08.09.2026}

\subsubsection{Sequences}

\begin{definition}
  Let \(X\) and \(Y\) be sets. A function \(f\colon X\to Y\) is a relation
  \(f\subseteq X\times Y\) such that, for every \(x\in X\), there exists
  exactly one \(y\in Y\) with \((x,y)\in f\). This unique \(y\) is denoted
  by \(f(x)\).

  The set \(X\) is the \emph{domain} of \(f\), and \(Y\) is its
  \emph{codomain}. The set
  \[
    f(X)=\{f(x):x\in X\}
  \]
  is its \emph{range} or \emph{image}.
\end{definition}

\begin{definition}
  A sequence of real numbers is a function \(x\colon\NN\to\RR\).
  For \(n\in\NN\), the value \(x(n)\), usually denoted by \(x_n\), is the
  \(n\)-th term of the sequence. We denote the sequence by
  \((x_n)_{n\in\NN}\), or simply by \((x_n)\).
\end{definition}

\subsubsection{Limit of a sequence}

\begin{definition}
  Let \((x_n)\) be a sequence of real numbers. We say that \((x_n)\)
  \emph{converges} to \(a\in\RR\) if, for every \(\epsilon>0\), there
  exists \(N\in\NN\) such that, for every \(n\in\NN\) with \(n > N\),
  \[
    \abs{x_n-a}<\epsilon.
  \]
  In this case, we write \(x_n\to a\) as \(n\to\infty\), or
  \(\lim_{n\to\infty}x_n=a\).
\end{definition}

Let us look at some examples.

\begin{example}
  Let \((x_n)\) be defined by
  \[
    x_n = \frac{2n}{n + 1}.
  \]
  We will show that
  \[
    \lim_{n \to \infty} x_n = 2.
  \]
  Let \(\epsilon > 0\). Choose \(N \in \NN\) such that
  \(N > \frac{2}{\epsilon}\). Now for every \(n > N\), we have
  \[
    \begin{aligned}
      \abs{x_n - 2} &= \abs{\frac{2n}{n + 1} - 2} \\
      &= \frac{2}{n + 1} < \frac{2}{n} < \frac{2}{\frac{2}{\epsilon}}
      = \epsilon.
    \end{aligned}
  \]
  Thus \(x_n \to 2\) when \(n \to \infty\).
\end{example}

\begin{example}
  Let \((x_n)\) be defined by
  \[
    \begin{aligned}
      x_0 &= \cdots = x_7 = 42, \\
      x_n &= 7 + \frac{1}{n}, \quad n \geq 8.
    \end{aligned}
  \]
  We will show that
  \[
    \lim_{n \to \infty} x_n = 7.
  \]
  Let \(\epsilon > 0\). Choose \(N \in \NN\) such that
  \(N > \maxset{8, \frac{1}{\epsilon}}\). Now for every \(n > N\), we have
  \[
    \abs{x_n - 7} = \frac{1}{n} < \epsilon.
  \]
  Thus \(x_n \to 7\) as \(n \to \infty\).
\end{example}

Let us now tackle a simple divergence proof.

\begin{example}
  Let us show that the sequence \(x_n\) defined by \(x_n = n\) diverges.
  Assume that the sequence converged to some \(a \in \RR\).
  Taking \(\epsilon = 1\), there exists \(N \in \NN\) such that,
  for every \(n > N\),
  \[
    \abs{x_n - a} < \epsilon = 1.
  \]
  Choose \(n \in \NN\) such that
  \[
    n > \maxset{N, \abs{a} + 1}.
  \]
  By the triangle inequality,
  \[
    n=\abs{n}=\abs{(n-a)+a}\leq\abs{n-a}+\abs{a}.
  \]
  Subtracting \(\abs{a}\) from both sides and using
  \(n>\abs{a}+1\), we obtain
  \[
    \abs{x_n-a}=\abs{n-a}\geq n-\abs{a}>1,
  \]
  contradicting the assumption \(\abs{x_n - a} < 1\). Therefore
  \((x_n)\) diverges.
\end{example}

\begin{example}
  Let \((x_n)\) be defined by
  \[
    \begin{aligned}
      x_0 &= \cdots = x_{42} = 1, \\
      x_n &= \frac{n}{n - 42}, \quad n \geq 43.
    \end{aligned}
  \]
  We will show that
  \[
    \lim_{n \to \infty} x_n = 1.
  \]
  Let \(\epsilon > 0\). Choose \(N \in \NN\) such that
  \[
    N > \maxset{84, \frac{84}{\epsilon}}.
  \]
  Now, for every \(n>N\), we have \(n>84\), and hence
  \(\frac{n}{2}-42>0\). Therefore
  \[
    \begin{aligned}
      \abs{x_n - 1} &= \abs*{\frac{n}{n - 42} - 1} \\
      &= \frac{42}{n - 42} \\
      &= \frac{42}{\frac{n}{2} + \frac{n}{2} - 42} \\
      &< \frac{84}{n} < \epsilon.
    \end{aligned}
  \]
  Thus \(x_n \to 1\) as \(n \to \infty\).
\end{example}

\subsection{10.09.2026}

\subsubsection{Uniqueness of limits}

\begin{theorem}[Uniqueness of limits]
  A convergent sequence of real numbers has a unique limit. That is, if
  \(x_n \to a\) and \(x_n \to b\) as \(n \to \infty\), then \(a=b\).
\end{theorem}

\begin{proof}
  Suppose for contradiction that \(a\neq b\), and set
  \[
    \epsilon=\frac{\abs{a-b}}{2}>0.
  \]
  Since \(x_n\to a\) and \(x_n\to b\), there exist \(N_a,N_b\in\NN\) such
  that
  \[
    n>N_a \implies \abs{x_n-a}<\epsilon
  \]
  and
  \[
    n>N_b \implies \abs{x_n-b}<\epsilon.
  \]
  Choose \(n>\maxset{N_a,N_b}\). By the triangle inequality,
  \[
    \abs{a-b}
    \leq \abs{a-x_n}+\abs{x_n-b}
    < 2\epsilon
    =\abs{a-b},
  \]
  which is impossible. Therefore \(a=b\).
\end{proof}

\begin{theorem}[Convergent sequences are bounded]
  Every convergent sequence of real numbers is bounded.
\end{theorem}

\begin{proof}
  Suppose that \(x_n\to a\). Taking \(\epsilon=1\), there exists
  \(N\in\NN\) such that
  \[
    n>N \implies \abs{x_n-a}<1.
  \]
  By the triangle inequality, for every \(n>N\),
  \[
    \abs{x_n}
    \leq \abs{x_n-a}+\abs{a}
    <1+\abs{a}.
  \]
  Let
  \[
    M=\maxset{\abs{x_0},\abs{x_1},\ldots,\abs{x_N},1+\abs{a}}.
  \]
  Then \(\abs{x_n}\leq M\) for every \(n\in\NN\), so \((x_n)\) is
  bounded.
\end{proof}

\subsubsection{Algebra of limits}

\begin{theorem}[Algebraic limit laws]
  Suppose that \(x_n\to a\) and \(y_n\to b\), and let \(c\in\RR\). Then
  \begin{enumerate}
    \item \(x_n\pm y_n\to a\pm b\).
    \item \(cx_n\to ca\).
    \item \(x_ny_n\to ab\).
    \item If \(b\neq0\), then \(y_n\neq0\) for all sufficiently large
      \(n\), and
      \[
        \frac{x_n}{y_n}\to\frac{a}{b}.
      \]
  \end{enumerate}
\end{theorem}

\begin{proof}
  \begin{enumerate}
    \item
      Let \(\epsilon>0\). Since \(x_n\to a\) and \(y_n\to b\), there exist
      \(N_x,N_y\in\NN\) such that
      \[
        n>N_x\implies\abs{x_n-a}<\frac{\epsilon}{2}
      \]
      and
      \[
        n>N_y\implies\abs{y_n-b}<\frac{\epsilon}{2}.
      \]
      Thus, for \(n>\maxset{N_x,N_y}\),
      \[
        \abs{(x_n+y_n)-(a+b)}
        \leq\abs{x_n-a}+\abs{y_n-b}
        <\epsilon.
      \]
      Hence \(x_n+y_n\to a+b\). Replacing \(y_n\) and \(b\) by \(-y_n\)
      and \(-b\) gives \(x_n-y_n\to a-b\).

    \item
      If \(c=0\), then \(cx_n=ca=0\) for every \(n\). If \(c\neq0\), choose
      \(N\in\NN\) such that
      \[
        n>N\implies\abs{x_n-a}<\frac{\epsilon}{\abs{c}}.
      \]
      Then, for every \(n>N\),
      \[
        \abs{cx_n-ca}=\abs{c}\abs{x_n-a}<\epsilon.
      \]
      Therefore \(cx_n\to ca\).

    \item
      Since \((y_n)\) converges, it is bounded. Choose \(M>0\) such that
      \(\abs{y_n}\leq M\) for every \(n\). There exist \(N_1,N_2\in\NN\)
      such that
      \[
        n>N_1\implies\abs{x_n-a}<\frac{\epsilon}{2M}
      \]
      and
      \[
        n>N_2\implies\abs{y_n-b}<\frac{\epsilon}{2(\abs{a}+1)}.
      \]
      For \(n>\maxset{N_1,N_2}\),
      \begin{align*}
        \abs{x_ny_n-ab}
        &\leq\abs{x_ny_n-ay_n}+\abs{ay_n-ab} \\
        &=\abs{y_n}\abs{x_n-a}+\abs{a}\abs{y_n-b} \\
        &<M\frac{\epsilon}{2M}
        +\abs{a}\frac{\epsilon}{2(\abs{a}+1)} \\
        &<\epsilon.
      \end{align*}
      Therefore \(x_ny_n\to ab\).

    \item
      Suppose that \(b\neq0\). Since \(y_n\to b\), there exists
      \(N_0\in\NN\) such that
      \[
        n>N_0\implies\abs{y_n-b}<\frac{\abs{b}}{2}.
      \]
      The reverse triangle inequality then gives
      \[
        \abs{y_n}\geq\abs{b}-\abs{y_n-b}>\frac{\abs{b}}{2}>0
      \]
      for every \(n>N_0\). In particular, \(y_n\neq0\) eventually. Choose
      \(N_1\in\NN\) such that
      \[
        n>N_1\implies
        \abs{y_n-b}<\frac{\epsilon\abs{b}^2}{2}.
      \]
      For \(n>\maxset{N_0,N_1}\),
      \[
        \abs*{\frac{1}{y_n}-\frac{1}{b}}
        =\frac{\abs{y_n-b}}{\abs{y_n}\abs{b}}
        <\frac{2}{\abs{b}^2}\abs{y_n-b}
        <\epsilon.
      \]
      Hence \(1/y_n\to1/b\); the finitely many terms with \(n\leq N_0\) may be
      defined arbitrarily because they do not affect the limit. Applying the
      product law to \((x_n)\) and \((1/y_n)\) yields
      \[
        \frac{x_n}{y_n}\to\frac{a}{b}.
      \]
  \end{enumerate}
\end{proof}

\subsubsection{Limits and order}

\begin{theorem}
  Suppose that \(x_n\to x\). Then
  \begin{enumerate}
    \item if \(x_n\leq M\) for every \(n\in\NN\), then \(x\leq M\);
    \item if \(x_n\geq m\) for every \(n\in\NN\), then \(x\geq m\).
  \end{enumerate}
\end{theorem}

\begin{proof}
  \begin{enumerate}
    \item
      Suppose for contradiction that \(x>M\), and set
      \[
        \epsilon=\frac{x-M}{2}>0.
      \]
      Since \(x_n\to x\), there exists \(N\in\NN\) such that, for every
      \(n>N\),
      \[
        \abs{x_n-x}<\epsilon.
      \]
      It follows that
      \[
        x_n>x-\epsilon
        =x-\frac{x-M}{2}
        =\frac{x+M}{2}
        >M,
      \]
      contradicting \(x_n\leq M\). Therefore \(x\leq M\).

    \item
      Suppose for contradiction that \(x<m\), and set
      \[
        \epsilon=\frac{m-x}{2}>0.
      \]
      Since \(x_n\to x\), there exists \(N\in\NN\) such that, for every
      \(n>N\),
      \[
        \abs{x_n-x}<\epsilon.
      \]
      It follows that
      \[
        x_n<x+\epsilon
        =x+\frac{m-x}{2}
        =\frac{x+m}{2}
        <m,
      \]
      contradicting \(x_n\geq m\). Therefore \(x\geq m\).
  \end{enumerate}
\end{proof}

\begin{theorem}[Squeeze theorem]
  Suppose that
  \[
    x_n\leq y_n\leq z_n
  \]
  for every \(n\in\NN\). If \(x_n\to a\) and \(z_n\to a\), then
  \(y_n\to a\).
\end{theorem}

\begin{proof}
  Let \(\epsilon>0\). Since \(x_n\to a\), there exists \(N_x\in\NN\) such
  that
  \[
    n>N_x\implies a-\epsilon<x_n<a+\epsilon.
  \]
  Similarly, since \(z_n\to a\), there exists \(N_z\in\NN\) such that
  \[
    n>N_z\implies a-\epsilon<z_n<a+\epsilon.
  \]
  Therefore, for every \(n>\maxset{N_x,N_z}\),
  \[
    a-\epsilon<x_n\leq y_n\leq z_n<a+\epsilon.
  \]
  Hence \(\abs{y_n-a}<\epsilon\), so \(y_n\to a\).
\end{proof}

The inequalities in the squeeze theorem need only hold for all sufficiently
large \(n\), since changing finitely many terms does not affect convergence.

\begin{example}
  Consider the oscillating sequence
  \[
    x_n=\frac{(-1)^n}{n}, \qquad n\geq1.
  \]
  Its sign alternates, but
  \[
    -\frac{1}{n}\leq\frac{(-1)^n}{n}\leq\frac{1}{n}.
  \]
  Both bounding sequences converge to \(0\). Indeed, given \(\epsilon>0\),
  choosing \(N>1/\epsilon\) gives \(1/n<\epsilon\) whenever \(n>N\).
  The squeeze theorem therefore gives
  \[
    \lim_{n\to\infty}\frac{(-1)^n}{n}=0.
  \]
  Notice that the squeeze theorem lets us prove convergence without trying
  to handle even and odd indices separately.
\end{example}

\begin{example}
  Consider
  \[
    x_n=\sqrt{n^2+n}-n, \qquad n\geq1.
  \]
  Rationalising the expression gives
  \[
    x_n=\frac{n}{\sqrt{n^2+n}+n}.
  \]
  Since
  \[
    n<\sqrt{n^2+n}<n+1,
  \]
  we obtain
  \[
    \frac{n}{2n+1}<x_n<\frac{1}{2}.
  \]
  By the algebraic limit laws,
  \[
    \lim_{n\to\infty}\frac{n}{2n+1}
    =\lim_{n\to\infty}\frac{1}{2+1/n}
    =\frac{1}{2}.
  \]
  The upper bound is the constant sequence \(1/2\), so both bounds converge
  to \(1/2\). Hence the squeeze theorem yields
  \[
    \lim_{n\to\infty}\left(\sqrt{n^2+n}-n\right)=\frac{1}{2}.
  \]
\end{example}

\subsubsection{Subsequences}

\begin{definition}
  Let \((x_n)\) be a sequence, and let \((n_k)\) be a strictly increasing
  sequence of natural numbers:
  \[
    n_0<n_1<n_2<\cdots.
  \]
  The sequence \((x_{n_k})\), obtained by selecting the terms of \((x_n)\)
  whose indices are \(n_0,n_1,n_2,\ldots\), is called a \emph{subsequence}
  of \((x_n)\).
\end{definition}

\begin{theorem}
  If \(x_n\to x\), then every subsequence of \((x_n)\) converges to \(x\).
\end{theorem}

\begin{proof}
  Let \((x_{n_k})\) be a subsequence of \((x_n)\), and let \(\epsilon>0\).
  Since \(x_n\to x\), there exists \(N\in\NN\) such that
  \[
    n>N\implies\abs{x_n-x}<\epsilon.
  \]
  Because \((n_k)\) is strictly increasing, \(n_k\geq k\) for every
  \(k\in\NN\). Thus, whenever \(k>N\), we also have \(n_k>N\), and hence
  \[
    \abs{x_{n_k}-x}<\epsilon.
  \]
  Therefore \(x_{n_k}\to x\). Since the subsequence was arbitrary, every
  subsequence of \((x_n)\) converges to the same limit \(x\).
\end{proof}

\subsection{15.09.2026}

\subsubsection{Upper and lower bounds}

\begin{definition}
  Let \(\varnothing\neq E\subseteq\RR\). A number \(s\in\RR\) is an
  \emph{upper bound} of \(E\) if
  \[
    x\leq s \qquad\text{for every }x\in E.
  \]
  If such a number exists, \(E\) is \emph{bounded above}.
\end{definition}

\begin{definition}
  Let \(\varnothing\neq E\subseteq\RR\). A number \(p\in\RR\) is a
  \emph{lower bound} of \(E\) if
  \[
    p\leq x \qquad\text{for every }x\in E.
  \]
  If such a number exists, \(E\) is \emph{bounded below}.
\end{definition}

\begin{definition}
  Let \(\varnothing\neq E\subseteq\RR\) be bounded above. An upper bound
  \(s\) of \(E\) is the \emph{least upper bound}, or \emph{supremum}, of
  \(E\) if \(s\leq u\) for every upper bound \(u\) of \(E\). We write
  \[
    s=\sup E.
  \]
\end{definition}

\begin{definition}
  Let \(\varnothing\neq E\subseteq\RR\) be bounded below. A lower bound
  \(p\) of \(E\) is the \emph{greatest lower bound}, or \emph{infimum}, of
  \(E\) if \(p\geq\ell\) for every lower bound \(\ell\) of \(E\). We write
  \[
    p=\inf E.
  \]
\end{definition}

If \(E\) is not bounded above, we write \(\sup E=+\infty\). Similarly,
if \(E\) is not bounded below, we write \(\inf E=-\infty\).

\subsubsection{The completeness axiom}

\paragraph{Completeness axiom.}
Every nonempty subset of \(\RR\) that is bounded above has a supremum
in \(\RR\).

\begin{corollary}
  Every nonempty subset \(E\subseteq\RR\) that is bounded below has an
  infimum in \(\RR\).
\end{corollary}

\begin{proof}
  Let \(c\) be a lower bound of \(E\), and define
  \[
    F=\set{-x:x\in E}.
  \]
  Since \(c\leq x\) for every \(x\in E\), we have \(-x\leq-c\).
  Thus \(F\) is bounded above. It is also nonempty, so the completeness
  axiom gives \(\sup F\in\RR\). For every \(x\in E\),
  \[
    -x\leq\sup F\leq-c,
    \qquad\text{hence}\qquad
    x\geq-\sup F\geq c.
  \]
  Therefore \(-\sup F\) is a lower bound of \(E\). Since the argument
  applies to every lower bound \(c\) of \(E\), it is the greatest lower
  bound. Consequently,
  \[
    \inf E=-\sup F\in\RR.
  \]
\end{proof}

\subsubsection{Maximum and properties of the supremum}

\begin{definition}
  A number \(s\) is the \emph{greatest element}, or \emph{maximum}, of
  a set \(E\) if \(s\in E\) and \(x\leq s\) for every \(x\in E\).
  We write \(s=\max E\).
\end{definition}

\begin{theorem}
  Let \(\varnothing\neq E\subseteq\RR\) be bounded above. If
  \(s=\max E\), then \(s=\sup E\).
\end{theorem}

\begin{proof}
  Since \(x\leq s\) for every \(x\in E\), the number \(s\) is an upper
  bound of \(E\). Hence \(\sup E\leq s\). On the other hand, \(s\in E\),
  so \(s\leq\sup E\). Thus \(s=\sup E\).
\end{proof}

\begin{theorem}
  The supremum of a nonempty subset of \(\RR\) that is bounded above
  is unique.
\end{theorem}

\begin{proof}
  Suppose that \(s\) and \(s'\) are both least upper bounds of \(E\).
  Since \(s'\) is an upper bound and \(s\) is a least upper bound,
  \(s\leq s'\). Similarly, since \(s\) is an upper bound and \(s'\) is a
  least upper bound, \(s'\leq s\). Therefore \(s=s'\).
\end{proof}

\begin{theorem}
  Let \(\varnothing\neq E\subseteq\RR\) be bounded above. For every
  \(\epsilon>0\), there exists \(x\in E\) such that
  \[
    \sup E-\epsilon<x\leq\sup E.
  \]
\end{theorem}

\begin{proof}
  Since \(\sup E\) is an upper bound, \(x\leq\sup E\) for every
  \(x\in E\). Suppose, for a contradiction, that for some
  \(\epsilon>0\),
  \[
    x\leq\sup E-\epsilon \qquad\text{for every }x\in E.
  \]
  Then \(\sup E-\epsilon\) is an upper bound of \(E\), but
  \(\sup E-\epsilon<\sup E\). This contradicts the fact that \(\sup E\)
  is the least upper bound. Hence the claimed element \(x\) exists.
\end{proof}

\subsubsection{Monotone sequences}

\begin{definition}
  A real sequence \((x_n)\) is
  \begin{itemize}
    \item \emph{increasing} (nondecreasing) if \(x_n\leq x_{n+1}\) for
      every \(n\in\NN\);
    \item \emph{decreasing} (nonincreasing) if \(x_{n+1}\leq x_n\) for
      every \(n\in\NN\);
    \item \emph{strictly increasing} if \(x_n<x_{n+1}\) for every
      \(n\in\NN\);
    \item \emph{strictly decreasing} if \(x_{n+1}<x_n\) for every
      \(n\in\NN\).
  \end{itemize}
  A sequence is \emph{monotone} if it is increasing or decreasing, and
  \emph{strictly monotone} if it is strictly increasing or strictly decreasing.
\end{definition}

\subsection{17.09.2026}

\subsubsection{Monotone convergence}

\begin{theorem}[Monotone convergence theorem]\label{thm:monotone-convergence}
  Let \((x_n)\) be an increasing sequence of real numbers. If there is
  \(M\in\RR\) such that \(x_n\leq M\) for every \(n\in\NN\), then
  \((x_n)\) converges, and
  \[
    \lim_{n\to\infty}x_n=\sup\set{x_n:n\in\NN}\leq M.
  \]
\end{theorem}

\begin{proof}
  The set \(E=\set{x_n:n\in\NN}\) is nonempty and bounded above by
  \(M\). By completeness, \(s=\sup E\) exists in \(\RR\), and
  \(s\leq M\). Let \(\epsilon>0\). By the approximation property of
  the supremum, there is an index \(N\in\NN\) such that
  \[
    s-\epsilon<x_N\leq s.
  \]
  Since \((x_n)\) is increasing, for every \(n>N\),
  \[
    s-\epsilon<x_N\leq x_n\leq s.
  \]
  Hence \(\abs{x_n-s}=s-x_n<\epsilon\), proving \(x_n\to s\).
\end{proof}

\begin{corollary}[Decreasing case]\label{cor:decreasing-convergence}
  Let \((x_n)\) be a decreasing sequence of real numbers. If there is
  \(m\in\RR\) such that \(m\leq x_n\) for every \(n\in\NN\), then
  \((x_n)\) converges, and
  \[
    \lim_{n\to\infty}x_n=\inf\set{x_n:n\in\NN}\geq m.
  \]
  Consequently, every bounded monotone real sequence converges.
\end{corollary}

\begin{proof}
  The sequence \((-x_n)\) is increasing and bounded above by \(-m\).
  By \cref{thm:monotone-convergence},
  \[
    -x_n\longrightarrow\sup\set{-x_n:n\in\NN}.
  \]
  Multiplying by \(-1\) and using the relation between infima and
  suprema gives
  \[
    x_n\longrightarrow-\sup\set{-x_n:n\in\NN}
    =\inf\set{x_n:n\in\NN}\geq m.
  \]
  A bounded monotone sequence satisfies the hypotheses of either the
  increasing or the decreasing case, so it converges.
\end{proof}

\subsubsection{Examples of decreasing sequences}

In the following two examples the indices start at \(n=1\).

\begin{example}
  Let
  \[
    x_1=3,\qquad x_n=\frac{n}{n-1}\quad(n\geq2).
  \]
  Show that \((x_n)\) is strictly decreasing and find its limit.
\end{example}

\begin{solution}
  First, \(x_2=2<3=x_1\). For \(n\geq2\), all terms are positive and
  \[
    \frac{x_{n+1}}{x_n}
    =\frac{(n+1)(n-1)}{n^2}
    =1-\frac{1}{n^2}<1.
  \]
  Thus \(x_{n+1}<x_n\) for every \(n\geq1\). Also \(x_1>1\), and
  \[
    x_n=1+\frac{1}{n-1}>1\qquad(n\geq2).
  \]
  Hence the sequence is bounded below and converges by
  \cref{cor:decreasing-convergence}. We identify its limit in two ways.

  \paragraph{Method 1: the squeeze theorem.}
  For \(n\geq2\), we have \(n-1\geq n/2\), so
  \[
    1<x_n=1+\frac{1}{n-1}\leq1+\frac{2}{n}.
  \]
  Both bounding sequences tend to \(1\), so \(x_n\to1\).

  \paragraph{Method 2: the infimum.}
  The number \(1\) is a lower bound of \(E=\set{x_n:n\geq1}\).
  If \(a>1\), choose an integer \(n>1+1/(a-1)\). Then
  \[
    x_n<a
    \iff \frac{1}{n-1}<a-1
    \iff n>1+\frac{1}{a-1}.
  \]
  Thus no \(a>1\) is a lower bound of \(E\). Therefore \(\inf E=1\),
  and the decreasing case of the monotone convergence theorem yields
  \(\lim_{n\to\infty}x_n=1\).
\end{solution}

\begin{example}
  Let
  \[
    x_n=\frac{(n+1)^2}{n^2+1},\qquad n\geq1.
  \]
  Show that \((x_n)\) is decreasing and determine its limit.
\end{example}

\begin{solution}
  Since every term is positive, we can compare successive terms using
  their ratio:
  \begin{align*}
    \frac{x_{n+1}}{x_n}
    &=\frac{(n+2)^2(n^2+1)}{((n+1)^2+1)(n+1)^2}\\
    &=\frac{n^4+4n^3+5n^2+4n+4}
    {n^4+4n^3+7n^2+6n+2}.
  \end{align*}
  The denominator exceeds the numerator by
  \[
    2n^2+2n-2=2(n^2+n-1)>0\qquad(n\geq1).
  \]
  Therefore \(x_{n+1}/x_n<1\), so the sequence is strictly decreasing.
  Moreover,
  \[
    x_n=1+\frac{2n}{n^2+1}>1,
  \]
  so it is bounded below and converges to
  \(\inf\set{x_n:n\geq1}\).

  We now show that this infimum is \(1\). It is already a lower bound.
  Given \(a>1\), choose an integer \(n>\maxset{1,2/(a-1)}\). Then
  \[
    x_n=1+\frac{2n}{n^2+1}<1+\frac{2}{n}<a.
  \]
  Hence \(a\) is not a lower bound, and consequently
  \[
    \lim_{n\to\infty}x_n=\inf\set{x_n:n\geq1}=1.
  \]
  Alternatively, \(1<x_n<1+2/n\) gives the same limit directly by the
  squeeze theorem.
\end{solution}

\begin{remark}
  The restriction \(n\geq1\) matters for monotonicity in the second
  example. If we also define \(x_0\) by the same formula, then
  \(x_0=1<x_1=2\), so the full sequence is not decreasing. Its tail
  starting at \(n=1\) is strictly decreasing, and its limit is still \(1\).
\end{remark}

\subsubsection{The Bolzano--Weierstrass theorem}

\begin{theorem}[Bolzano--Weierstrass]
  Every bounded sequence of real numbers has a convergent subsequence.
\end{theorem}

\begin{proof}
  Let \((x_n)\) be bounded. Choose \(m,M\in\RR\) such that
  \(m\leq x_n\leq M\) for every \(n\in\NN\), and define the
  \emph{tail suprema}
  \[
    a_k=\sup\set{x_n:n\geq k},\qquad k\in\NN.
  \]
  Each tail is nonempty and bounded, so completeness ensures that
  \(a_k\) exists and \(m\leq a_k\leq M\). Since
  \[
    \set{x_n:n\geq k+1}\subseteq\set{x_n:n\geq k},
  \]
  the number \(a_k\) is an upper bound of the smaller tail as well.
  Therefore \(a_{k+1}\leq a_k\). The sequence \((a_k)\) is decreasing
  and bounded below, so \cref{cor:decreasing-convergence} gives a limit
  \[
    a=\lim_{k\to\infty}a_k\in\RR.
  \]

  We construct a subsequence \((x_{n_j})_{j\geq1}\) recursively.
  By the approximation property of the supremum, choose \(n_1\geq1\)
  such that
  \[
    a_1-1<x_{n_1}\leq a_1.
  \]
  Once \(n_j\) has been chosen, apply the same property to the tail
  starting at \(n_j+1\). There is an index \(n_{j+1}\geq n_j+1\)
  such that
  \[
    a_{n_j+1}-\frac{1}{j+1}<x_{n_{j+1}}\leq a_{n_j+1}.
  \]
  Thus the indices are strictly increasing, as required for a subsequence.

  To verify convergence, observe that
  \(a_{n_{j+1}}\leq a_{n_j+1}\) and
  \(x_{n_{j+1}}\leq a_{n_{j+1}}\). Consequently, for every \(j\geq1\),
  \[
    a_{n_j}-\frac{1}{j}<x_{n_j}\leq a_{n_j}.
  \]
  For \(j=1\), the lower inequality follows from
  \(a_{n_1}\leq a_1\); for later indices it follows from the recursive
  choice above. Since \((a_{n_j})\) is a subsequence of the convergent
  sequence \((a_k)\), it tends to \(a\). Also \(1/j\to0\), so both
  bounding sequences tend to \(a\). The squeeze theorem now yields
  \(x_{n_j}\to a\), proving the claim.
\end{proof}

\subsection{22.09.2026}

\subsubsection{Bernoulli's inequality}

\begin{theorem}[Bernoulli's inequality]\label{thm:bernoulli}
  For every \(x\in\RR\) with \(x>-1\) and every \(n\in\NN\),
  \[
    (1+x)^n\geq1+nx.
  \]
\end{theorem}

\begin{proof}
  Fix \(x>-1\). The cases \(n=0\) and \(n=1\) are equalities.
  Suppose that \((1+x)^n\geq1+nx\) for some \(n\geq1\).
  Since \(1+x>0\), multiplication preserves the inequality, and
  \begin{align*}
    (1+x)^{n+1}
    &=(1+x)(1+x)^n\\
    &\geq(1+x)(1+nx)\\
    &=1+(n+1)x+nx^2\\
    &\geq1+(n+1)x.
  \end{align*}
  The result follows by induction. If \(x\neq0\), then \(nx^2>0\)
  in the induction step, so the inequality is strict for \(n\geq2\).
\end{proof}

\subsubsection{A sequence defining the number \texorpdfstring{\(e\)}{e}}

Consider the sequence
\[
  x_n=\left(1+\frac{1}{n}\right)^n,\qquad n\geq1.
\]

\begin{lemma}
  The sequence \((x_n)\) is strictly increasing.
\end{lemma}

\begin{proof}
  All terms are positive. For \(n\geq1\),
  \begin{align*}
    \frac{x_{n+1}}{x_n}
    &=\frac{\left(1+\frac{1}{n+1}\right)^{n+1}}
    {\left(1+\frac{1}{n}\right)^n}\\
    &=\left(\frac{1+\frac{1}{n+1}}{1+\frac{1}{n}}\right)^{n+1}
    \left(1+\frac{1}{n}\right)\\
    &=\left(\frac{n(n+2)}{(n+1)^2}\right)^{n+1}
    \frac{n+1}{n}\\
    &=\left(1-\frac{1}{(n+1)^2}\right)^{n+1}\frac{n+1}{n}\\
    &>\left(1-\frac{n+1}{(n+1)^2}\right)\frac{n+1}{n}
    =1.
  \end{align*}
  The strict inequality follows from \cref{thm:bernoulli}, since
  \(-1<-1/(n+1)^2<0\) and \(n+1\geq2\).
  Hence \(x_{n+1}>x_n\).
\end{proof}

\begin{lemma}
  The sequence \((x_n)\) is bounded above.
\end{lemma}

\begin{proof}
  We use the auxiliary sequence
  \[
    y_n=\left(1-\frac{1}{n}\right)^n,\qquad n\geq2.
  \]
  This sequence is positive and strictly increasing. Indeed, by
  \cref{thm:bernoulli}, for \(n\geq2\),
  \begin{align*}
    \frac{y_{n+1}}{y_n}
    &=\left(\frac{n^2}{n^2-1}\right)^n\frac{n}{n+1}\\
    &=\left(1+\frac{1}{n^2-1}\right)^n\frac{n}{n+1}\\
    &\geq\left(1+\frac{n}{n^2-1}\right)\frac{n}{n+1}\\
    &=1+\frac{1}{(n+1)(n^2-1)}>1.
  \end{align*}
  Consequently, \((1/y_n)\) is decreasing. Also, for \(n\geq2\),
  \[
    \left(1+\frac{1}{n}\right)\left(1-\frac{1}{n}\right)
    =1-\frac{1}{n^2}<1.
  \]
  Dividing by \(1-1/n>0\) and raising to the power \(n\) gives
  \[
    x_n<\left(1-\frac{1}{n}\right)^{-n}
    =\frac{1}{y_n}\leq\frac{1}{y_2}
    =\left(1-\frac{1}{2}\right)^{-2}=4.
  \]
  Finally, \(x_1=2<4\), so the bound holds for every \(n\geq1\).
\end{proof}

By \cref{thm:monotone-convergence}, the increasing, bounded sequence
\((x_n)\) converges to a real number.

\begin{definition}
  The limit
  \[
    e=\lim_{n\to\infty}\left(1+\frac{1}{n}\right)^n
  \]
  is called \emph{Euler's number} (\emph{Neperin luku} in Finnish).
\end{definition}

\subsection{24.09.2026}

\subsubsection{Sequences tending to infinity}

\begin{definition}
  A real sequence \((x_n)\) \emph{tends to positive infinity} if, for
  every \(M\in\RR\), there exists \(N\in\NN\) such that
  \[
    x_n>M\qquad\text{whenever }n>N.
  \]
  We write \(x_n\to+\infty\), or
  \(\lim_{n\to\infty}x_n=+\infty\).
\end{definition}

\begin{definition}
  A real sequence \((x_n)\) \emph{tends to negative infinity} if, for
  every \(M\in\RR\), there exists \(N\in\NN\) such that
  \[
    x_n<M\qquad\text{whenever }n>N.
  \]
  We write \(x_n\to-\infty\), or
  \(\lim_{n\to\infty}x_n=-\infty\).
\end{definition}

\begin{remark}
  Neither definition requires the sequence to be monotone. The symbols
  \(+\infty\) and \(-\infty\) are not real numbers, so these statements
  do not mean convergence to a real limit.
\end{remark}

\begin{theorem}
  If \(x_n\to+\infty\) or \(x_n\to-\infty\), then \((x_n)\)
  diverges as a sequence of real numbers.
\end{theorem}

\begin{proof}
  Suppose first that \(x_n\to+\infty\), and assume for a contradiction
  that \(x_n\to x\in\RR\). Taking \(M=x+1\) in the definition
  of an infinite limit gives \(N_1\in\NN\) such that
  \(x_n>x+1\) whenever \(n>N_1\). Convergence to \(x\), with
  \(\varepsilon=1\), gives \(N_2\in\NN\) such that
  \(\abs{x_n-x}<1\) whenever \(n>N_2\). For
  \(n>\maxset{N_1,N_2}\), we would have both
  \(\abs{x_n-x}>1\) and \(\abs{x_n-x}<1\), a contradiction.

  If \(x_n\to-\infty\), the same argument uses \(M=x-1\).
  Eventually \(x_n<x-1\), hence \(\abs{x_n-x}>1\), again
  contradicting convergence to \(x\).
\end{proof}

\subsubsection{A necessary condition for convergence}

\begin{theorem}\label{thm:convergent-cauchy}
  Let \((x_n)\) be a convergent real sequence. For every
  \(\varepsilon>0\), there exists \(N\in\NN\) such that
  \[
    \abs{x_n-x_{n+k}}<\varepsilon
    \qquad\text{for all }n>N\text{ and }k\in\NN.
  \]
\end{theorem}

\begin{proof}
  Write \(x=\lim_{n\to\infty}x_n\), and let \(\varepsilon>0\).
  Choose \(N\in\NN\) such that \(\abs{x_m-x}<\varepsilon/2\)
  for every \(m>N\). If \(n>N\) and \(k\in\NN\), then
  \(n+k>N\), so the triangle inequality gives
  \[
    \abs{x_n-x_{n+k}}
    \leq\abs{x_n-x}+\abs{x_{n+k}-x}
    <\frac{\varepsilon}{2}+\frac{\varepsilon}{2}
    =\varepsilon.
  \]
\end{proof}

\begin{definition}
  A real sequence satisfying the condition in
  \cref{thm:convergent-cauchy} is called a \emph{Cauchy sequence}.
  Equivalently, for every \(\varepsilon>0\), there exists \(N\in\NN\)
  such that
  \[
    \abs{x_p-x_q}<\varepsilon\qquad\text{whenever }p,q>N.
  \]
  Indeed, take \(n=\min\set{p,q}\) and \(k=\abs{p-q}\) to obtain
  this form from the preceding one. Conversely, set \(p=n\) and
  \(q=n+k\).
\end{definition}

\subsubsection{Nested closed intervals}

\begin{lemma}[Nested interval lemma]\label{lem:nested-intervals}
  Let \([a_n,b_n]\), \(n\geq1\), be nonempty closed intervals such that
  \[
    [a_{n+1},b_{n+1}]\subseteq[a_n,b_n]\qquad(n\geq1),
    \qquad \lim_{n\to\infty}(b_n-a_n)=0.
  \]
  Then \(\bigcap_{n=1}^{\infty}[a_n,b_n]\) contains exactly one
  real number.
\end{lemma}

\begin{proof}
  The nesting gives
  \[
    a_n\leq a_{n+1}\leq b_{n+1}\leq b_n.
  \]
  Thus \((a_n)\) is increasing and \((b_n)\) is decreasing. Both
  sequences lie in \([a_1,b_1]\), so the monotone convergence theorem
  gives real limits
  \[
    a=\lim_{n\to\infty}a_n,
    \qquad b=\lim_{n\to\infty}b_n.
  \]
  Since \(a_n\leq b_n\), preservation of inequalities under limits
  yields \(a\leq b\). Monotonicity also gives
  \(a_n\leq a\leq b\leq b_n\) for every \(n\), hence
  \[
    0\leq b-a\leq b_n-a_n.
  \]
  The right-hand side tends to zero, so \(a=b\). In particular,
  \(a\in[a_n,b_n]\) for every \(n\), proving existence of a point
  in the intersection.

  For uniqueness, let \(y\neq a\). Since \(\abs{y-a}>0\) and
  \(b_n-a_n\to0\), choose \(n\) with
  \(b_n-a_n<\abs{y-a}\). If \(y\) also belonged to \([a_n,b_n]\),
  then \(a,y\in[a_n,b_n]\) would imply
  \(\abs{y-a}\leq b_n-a_n\), a contradiction. Therefore \(y\)
  does not belong to the intersection, and
  \[
    \bigcap_{n=1}^{\infty}[a_n,b_n]=\set{a}.
  \]
\end{proof}

\subsubsection{Cauchy's convergence criterion}

\begin{theorem}[Cauchy's convergence criterion]
  A real sequence \((x_n)\) converges if and only if it is a Cauchy
  sequence: for every \(\varepsilon>0\), there exists \(N\in\NN\)
  such that
  \[
    \abs{x_n-x_{n+k}}<\varepsilon
    \qquad\text{for all }n>N\text{ and }k\in\NN.
  \]
\end{theorem}

\begin{proof}
  Necessity was proved in \cref{thm:convergent-cauchy}. Conversely,
  suppose that \((x_n)\) satisfies the Cauchy condition.

  \paragraph{Step 1: boundedness.}
  Apply the condition with \(\varepsilon=1\), and choose a corresponding
  \(N_0\in\NN\). For every \(m\geq N_0+1\), take
  \(n=N_0+1\) and \(k=m-(N_0+1)\). Then
  \[
    \abs{x_m-x_{N_0+1}}<1,
    \qquad \abs{x_m}<\abs{x_{N_0+1}}+1.
  \]
  The finitely many earlier terms are bounded as well. In particular,
  \[
    C=\max\bigl\{\abs{x_0},\ldots,\abs{x_{N_0}},
    \abs{x_{N_0+1}}+1\bigr\}
  \]
  satisfies \(\abs{x_m}\leq C\) for every \(m\in\NN\).

  \paragraph{Step 2: nested intervals from the tails.}
  Define, for \(n\geq1\),
  \[
    p_n=\inf\set{x_m:m\geq n},
    \qquad s_n=\sup\set{x_m:m\geq n}.
  \]
  Each tail is nonempty and bounded, so its infimum and supremum
  exist in \(\RR\). As the tails shrink, their infima increase and
  their suprema decrease:
  \[
    p_n\leq p_{n+1}\leq s_{n+1}\leq s_n.
  \]
  Consequently, the nonempty closed intervals \([p_n,s_n]\) are nested.

  \paragraph{Step 3: the interval lengths tend to zero.}
  Let \(\varepsilon>0\). By the Cauchy condition, choose \(N\)
  so that \(\abs{x_r-x_t}<\varepsilon/3\) whenever \(r,t>N\).
  Fix \(n>N\). By the approximation properties of the supremum
  and infimum, there are indices \(r,t\geq n\) such that
  \[
    0\leq s_n-x_r<\frac{\varepsilon}{3},
    \qquad 0\leq x_t-p_n<\frac{\varepsilon}{3}.
  \]
  Since \(r,t>N\), we obtain
  \begin{align*}
    0\leq s_n-p_n
    &\leq\abs{s_n-x_r}+\abs{x_r-x_t}+\abs{x_t-p_n}\\
    &<\frac{\varepsilon}{3}+\frac{\varepsilon}{3}
    +\frac{\varepsilon}{3}=\varepsilon.
  \end{align*}
  Hence \(s_n-p_n\to0\).

  \paragraph{Step 4: convergence to the common point.}
  By \cref{lem:nested-intervals}, there is exactly one \(x\in\RR\)
  such that
  \[
    \bigcap_{n=1}^{\infty}[p_n,s_n]=\set{x}.
  \]
  Given \(\varepsilon>0\), choose \(n\) with
  \(s_n-p_n<\varepsilon\). For every \(m\geq n\), both
  \(x\) and \(x_m\) belong to \([p_n,s_n]\), so
  \[
    \abs{x_m-x}\leq s_n-p_n<\varepsilon.
  \]
  This proves \(x_m\to x\), completing the sufficiency direction.
\end{proof}

\end{document}
